\section{The Debye--Waller factor}\label{app:D.15}

A momentum transfer $\vect{q}$ to a nucleus acts through the operator $\ee^{\ii\vect{q}\cdot\vect{u}_{\vect{R}\kappa}}$ (\cref{app:C.6}). The average of the exponential of an operator is in general not the exponential of anything simple. For a harmonic crystal it is, by Bloch's identity \cref{eq:C.6-bloch}, $\langle\ee^{\op{A}}\rangle = \ee^{\langle\op{A}^2\rangle/2}$ for $\op{A}$ linear in the phonon operators. Here we prove the identity and derive the generating formula of \cref{sec:5.6} for two different nuclei.

\subsection*{Proof for one mode}

Take one oscillator, $H = \omega_0(d^\dagger d + \tfrac12)$ as in \cref{app:C.6}, in thermal equilibrium at temperature $T = 1/\beta$, with $n_{\mathrm{B}} = n_{\mathrm{B}}(\omega_0)$, and let $\op{A} = \eta_1d + \eta_2d^\dagger$ with arbitrary complex $\eta_1$ and $\eta_2$, so that $\op{A}$ need not be Hermitian. Consider $F(\theta) = \langle\ee^{\theta\op{A}}\rangle$ as a function of a real parameter $\theta$. The proof uses two facts. The commutators $[\op{A}, d] = -\eta_2$ and $[\op{A}, d^\dagger] = \eta_1$ are numbers. By \cref{eq:C.6-hadamard}, the transformation by $\ee^{\theta\op{A}}$ therefore only shifts the ladder operators,
\begin{equation*}
  d\,\ee^{\theta\op{A}} = \ee^{\theta\op{A}}\,(d + \theta\eta_2), \qquad d^\dagger\,\ee^{\theta\op{A}} = \ee^{\theta\op{A}}\,(d^\dagger - \theta\eta_1) .
\end{equation*}
The cyclic property of the trace, in turn, moves a ladder operator past the Boltzmann factor at a known cost, as in the proof of Wick's theorem in \cref{app:D.11}: $d\,\ee^{-\beta H} = \ee^{-\beta\omega_0}\ee^{-\beta H}d$, since removing a quantum lowers $H$ by $\omega_0$, and $d^\dagger\ee^{-\beta H} = \ee^{+\beta\omega_0}\ee^{-\beta H}d^\dagger$. As a result, $\langle\ee^{\theta\op{A}}d\rangle = \Tr[d\,\ee^{-\beta H}\ee^{\theta\op{A}}]/\mathcal{Z} = \ee^{-\beta\omega_0}\langle d\,\ee^{\theta\op{A}}\rangle$, and combined with the shift,
\begin{equation*}
  \langle d\,\ee^{\theta\op{A}}\rangle = \ee^{-\beta\omega_0}\langle d\,\ee^{\theta\op{A}}\rangle + \theta\eta_2F(\theta), \qquad
  \langle d\,\ee^{\theta\op{A}}\rangle = \frac{\theta\eta_2F(\theta)}{1 - \ee^{-\beta\omega_0}} = \theta\eta_2\,(n_{\mathrm{B}} + 1)\,F(\theta) .
\end{equation*}
The same steps for $d^\dagger$ give $\langle d^\dagger\ee^{\theta\op{A}}\rangle = \ee^{\beta\omega_0}\langle d^\dagger\ee^{\theta\op{A}}\rangle - \theta\eta_1F(\theta)$, and $\langle d^\dagger\ee^{\theta\op{A}}\rangle = \theta\eta_1\,n_{\mathrm{B}}\,F(\theta)$. The derivative of $F$ is then
\begin{equation*}
  F'(\theta) = \eta_1\langle d\,\ee^{\theta\op{A}}\rangle + \eta_2\langle d^\dagger\ee^{\theta\op{A}}\rangle = \theta\,\eta_1\eta_2\,(2n_{\mathrm{B}} + 1)\,F(\theta) = \theta\,\langle\op{A}^2\rangle\,F(\theta) ,
\end{equation*}
since $\langle\op{A}^2\rangle = \eta_1\eta_2\langle dd^\dagger + d^\dagger d\rangle = \eta_1\eta_2(2n_{\mathrm{B}} + 1)$, the averages of $d^2$ and $d^{\dagger2}$ being zero. With $F(0) = 1$ the solution is $F(\theta) = \ee^{\theta^2\langle\op{A}^2\rangle/2}$, and at $\theta = 1$
\begin{equation}
  \langle\ee^{\op{A}}\rangle = \ee^{\langle\op{A}^2\rangle/2} .
  \label{eq:D.15-bloch}
\end{equation}
For several independent modes in thermal equilibrium, $\op{A} = \sum_\gamma\op{A}_\gamma$ is a sum of commuting pieces, one per mode, the exponential factorizes, and the average in the Gibbs state of \cref{sec:2.5} factorizes too. The exponents add, and $\sum_\gamma\langle\op{A}_\gamma^2\rangle = \langle\op{A}^2\rangle$, because the cross terms $\langle\op{A}_\gamma\op{A}_{\gamma'}\rangle = \langle\op{A}_\gamma\rangle\langle\op{A}_{\gamma'}\rangle$ vanish for $\gamma \ne \gamma'$. \Cref{eq:D.15-bloch} therefore holds for any $\op{A}$ linear in the operators of a set of harmonic modes. Mermin gave a short evaluation of averages of this form by a different route~\cite{mermin1966short}.

The hypotheses are the ones the proof used. The operator must be linear in $d$ and $d^\dagger$, so that its commutators are numbers, and the state must be the thermal state of a quadratic Hamiltonian, so that the trace moves each ladder operator at the cost of a number. In an anharmonic crystal neither holds exactly. The logarithm of $\langle\ee^{\op{A}}\rangle$ is then a series in the cumulants of $\op{A}$, of which the harmonic crystal has only the second. The anharmonic terms of \cref{sec:5.7} make the fourth cumulant nonzero, and the Debye--Waller exponent acquires a correction that grows with temperature faster than $T$.

\subsection*{The generating formula}

For a crystal, the operator is $\op{A} = \ii\vect{q}\cdot\vect{u}_{\vect{R}\kappa}$ with the displacement of \cref{eq:A.6-u}. Its square averages to $-\vect{q}\cdot B_\kappa\cdot\vect{q}$, with the displacement tensor of \cref{eq:5.3-B}, and \cref{eq:D.15-bloch} gives the Debye--Waller factor \cref{eq:5.6-W}. The product formula \cref{eq:C.6-product}, which follows from \cref{eq:D.15-bloch} and the Baker--Campbell--Hausdorff formula, applies with $\op{A} = -\ii\vect{q}\cdot\vect{u}_{\vect{R}\kappa}(t)$ and $\op{B} = \ii\vect{q}\cdot\vect{u}_{\vect{0}\kappa'}(0)$, both linear. The averages of their squares are $-2W_\kappa(\vect{q})$ and $-2W_{\kappa'}(\vect{q})$, since the state is stationary, and the average of their product is $\vect{q}\cdot\mathcal{C}_{\kappa\kappa'}(\vect{R},t)\cdot\vect{q}$, with the displacement correlation of \cref{eq:5.6-C}. Hence
\begin{equation}
  \bigl\langle\ee^{-\ii\vect{q}\cdot\vect{u}_{\vect{R}\kappa}(t)}\,\ee^{\ii\vect{q}\cdot\vect{u}_{\vect{0}\kappa'}(0)}\bigr\rangle = \ee^{-W_\kappa(\vect{q}) - W_{\kappa'}(\vect{q})}\;\ee^{\vect{q}\cdot\mathcal{C}_{\kappa\kappa'}(\vect{R},t)\cdot\vect{q}} ,
  \label{eq:D.15-gen}
\end{equation}
and for one nucleus, $\vect{R} = 0$ and $\kappa' = \kappa$, this is the generating formula \cref{eq:5.6-gen}.
